\begin{theorem}[Four partner fluxes in the actual swept presentations]%
    \label{thm:peps_torus_swept_string_endpoint_holonomy}
    \lean{TNLean.PEPS.regularWalkHolonomy_torusSweptStringInitial,
        TNLean.PEPS.torusSweptPatchEndpoint_mem_sweptStringVertices,
        TNLean.PEPS.regularWalkHolonomy_torusSweptStringOperators}
    \leanok
    \uses{def:peps_torus_swept_string_assignment,
        def:peps_torus_swept_patch_endpoint_geometry,
        def:peps_torus_swept_patch_endpoint_walks,
        thm:peps_regular_walk_holonomy}
    Let $w\geq7$ and $H\geq6$. Complete the displayed patch by the
    four actual clockwise plaquette loops based at
    \begin{align}
        A_2 & =(x-1,y+1), & A_1 & =(x+3,y+1), \\
        B_1 & =(x+1,y+1), & B_2 & =(x+1,y-2).
    \end{align}
    For the literal initial assignment their holonomies, in this order,
    are $(g^{-1},g,h,h^{-1})$.
    Among these four base vertices, precisely $B_1$ lies in either
    swept set $S_1$ or $S_2$. For either swept assignment the actual
    clockwise holonomies are
    $(g^{-1},g,ghg^{-1},h^{-1})$.
    These are auxiliary finite-torus identities for the virtual
    string-deformation figure in \cite[Section~6.6]{Schuch2010PEPS},
    following the string-motion equation, lines 2361--2386.
    % Source label: eq:anyons:fluxon-braiding-lazy.
    They do not identify the prescribed physical braid.
\end{theorem}

\begin{proof}\leanok
    \uses{thm:peps_torus_swept_string_transport,
        thm:peps_zmod_small_difference,
        thm:peps_torus_swept_patch_endpoint_geometry,
        def:peps_torus_swept_patch_endpoint_walks,
        thm:peps_torus_plaquette_holonomy,thm:peps_regular_walk_holonomy}
    For a clockwise plaquette loop, multiply its four directed
    transports in reverse traversal order. Substituting the actual
    right and up tables yields the four initial values.
    The period bounds distinguish every required integer offset,
    proving that only $B_1$ belongs to either sweep.
    Closed-loop gauge covariance is
    $H(u^{(j)},p)=k_{\rm base}^{-1}H(u^{(0)},p)k_{\rm base}$.
    Hence only the loop at $B_1$ is conjugated, by $g$.
\end{proof}
